\scanpage{22}
\mspage{21}

question s'il en existe une semblerait \unsure{accessoire|nécessaire}. Aussi, finalement, j'ai envie de poser

\noindent\begin{tabular}{@{}||@{\hspace{.55em}}p{.93\linewidth}@{}}
\uline{A\deleted{C}\unsure{t|l|x} 4}\quad Si $\deleted{G}F,\deleted{F}G\in\unsure{\mathcal{F}|\mathscr{F}}$, alors $\deleted{F}G\ll\deleted{G}F$ sss pour tout $X\in\widetilde{\deleted{F}G}$, $\deleted{\unsure{\exists|s}}$ $\exists Y\in\widetilde F$ avec $X\ll Y$
\end{tabular}

Ainsi, la relation $\ll$ sur $\mathcal{F}$ se déduit de la relation $\ll$ sur $M$. Donc, si on \unsure{suppose|s'impose} les figures de type fini, $(\mathcal{F},\leq,\ll)$ (i.e. tout l'atelier) se réconstitue entièrement à partir de $(M,\leq,\ll,\mathrel{\tikz[baseline=-.55ex,x=1em,y=1em]{\draw[line width=.065em,line cap=round,line join=round] (0,.40)--(.48,0)--(0,-.35) (.51,.40)--(.03,.02)--(.49,-.35);}})$.
% slanted
\margin{\uline{Cor} $\forall F\in\mathcal{F},\ \mathrm{Omb}(F)=\bigl|\unsure{\mathrm{Multomb}|\mathrm{Multiomb}}\,F\bigr|$}

\noindent\begin{tabular}{@{}||@{\hspace{.55em}}p{.93\linewidth}@{}}
\uline{\deleted{C}At 5}\quad Soient $X\in M$, $F\in\mathcal{F}$, tels que $X\ll F$. Alors \deleted{dans} l'ensemble $F^X=\{Y\in\widetilde F\mid\underbrace{X\ll Y}_{\text{i.e. }X\in\mathrm{Omb}(Y)}\}$ (lequel est $\ne\varnothing$ par C4) admet un plus petit élément\unsure{,|.|*} (pour $\leq$)
\end{tabular}

En d'autres termes, l'ensemble de parties $\mathrm{Multom}\deleted{\unsure{b}}b(F)=\{\widetilde Y\mid Y\in\widetilde F\ \text{i.e.}\ Y\triangleleft F\}$ est une \uline{figure} dans $M$, \deleted{[\uline{NB} \unsure{étant entendu} que les $\widetilde Y^{\unsure{\circ|o}}=\widetilde Y\setminus\bigcup_{\unsure{X<|X\ll}}$} ou du moins, est \uuline{pré}\uline{admissible}. C'est en tous cas impliqué par \uline{C5}, et lui est équivalent si on admet que la relation (pour $Y,Y'\in\widetilde F$) $Y\leq Y'$, qui implique $Y\ll Y'$ \unsure{d'où|donc} $\deleted{\unsure{G|O}}\mathrm{Omb}(Y)\subset\mathrm{Omb}(Y')$ lui est en fait équivalente. Cela équivaut à la condition

\noindent\begin{tabular}{@{}||@{\hspace{.55em}}p{.93\linewidth}@{}}
\uline{\deleted{C}At 6}\quad Si $X,Y\in M$ et $X \mathrel{\tikz[baseline=-.55ex,x=1em,y=1em]{\draw[line width=.065em,line cap=round,line join=round] (0,.40)--(.48,0)--(0,-.35) (.51,.40)--(.03,.02)--(.49,-.35);}} Y$, alors $X\ll Y\Rightarrow X\leq Y$ (donc $X\ll Y\ \unsure{\Longleftrightarrow|\longleftrightarrow}\ X\leq Y$).

En d'autres termes,
\end{tabular}
